% META: {"id": "1NSI-ARCHOS-CO-008", "chapitre": "1NSI-ARCHITECTURE-OS", "type_objet": "correction", "exercice_ref": "1NSI-ARCHOS-EX-001", "capacites": ["1NSI-ARCHITECTURE-OS-C2"], "sources_inspiration": [], "status": "approved", "capacites_codes": ["C2"], "parcours": 2, "fichier_exercice": "chapitres/1NSI-ARCHITECTURE-OS/exercices/1NSI-ARCHOS-EX-008.tex"}
\begin{corrige}{1NSI-ARCHOS-EX-001}
\textbf{1.} Le \textbf{processeur} (unité centrale) exécute les instructions. La
\textbf{mémoire} stocke à la fois les instructions et les données.

\textbf{2.} Oui : avec $8$ cœurs, l'ordinateur est \textbf{multiprocesseur} et peut exécuter
plusieurs instructions en parallèle, un cœur par instruction (au sens large : plusieurs
programmes ou plusieurs parties d'un même programme peuvent progresser simultanément).

\textbf{3.} Instructions et données sont stockées sous la \textbf{même forme} (des suites de
bits) dans le même espace mémoire, sans marqueur qui les distinguerait intrinsèquement.
C'est le \textbf{contexte de lecture} (le séquenceur lit-il cette adresse pour l'exécuter
comme instruction, ou l'UAL la lit-elle comme donnée à calculer ?) qui détermine
l'interprétation de ces bits, pas leur contenu lui-même.
\end{corrige}
